\section{Background and Motivation}\label{sec:background}

% Compressed to ~1 page. Two subsections covering (a) what edge LLM
% inference looks like and where heterogeneity bites, and (b) why
% the obvious composition of prior systems doesn't compose.

\subsection{Pipeline-parallel inference on heterogeneous edge}
% Outline (~3 paragraphs):
% - Decoder block as the unit of partition; KV cache state per
%   (request, stage). Prefill is one large activation, decode is
%   per-token append (~2 KB at fp16 for OPT-350M's 1024-dim
%   hidden — the per-layer transfer unit).
% - Edge heterogeneity matters more than data-center heterogeneity:
%   AGX Orin (MAXN power mode) is 16x faster per layer than Jetson
%   Nano CUDA on OPT-350M decode (\S\ref{sec:eval:setup}). On a
%   homogeneous fleet placement is a rounding question; on
%   16x-different devices it is the only thing that matters.
% - Edge reliability is also worse: boards crash, SD cards fill,
%   heartbeat silence is common. Placement that ignores this
%   produces a stream-aborts-on-first-failure system.

\subsection{Why the obvious composition of prior systems fails}
% Outline (~2 paragraphs + 1-row comparison table):
% - EdgeShard's DP (latency-optimal placement, no R, assumes layer
%   redundancy) + Petals' swarm recovery (memory-redundant peers):
%   the 4 GB Nano can't hold Petals' redundancy.
% - Jupiter (throughput-optimal DP, no R) + a watchdog restart:
%   backup memory wasn't reserved at placement time, so on failure
%   either OOM or non-redundant degraded path.
% - The systemic point: $R$ and $\psi$ have *coupled* feasibility
%   regions on memory-tight edge; decoupling gives up feasibility
%   *or* performance, never both. This is the gap RADP fills.

\todo{Write the two subsections from outlines above; one short
table comparing EdgeShard / Petals / Jupiter / Helix / ours along
{placement DP, recovery model, runtime architecture, fleet target}
columns.}
